\section{Motivation and Problem Statement}
\label{sec:motivation}

\paragraph{What ``sub-threshold'' means.}
Throughout, \emph{sub-threshold} refers to the \emph{network-wide} corruption threshold, not a
committee's local one. A monolithic accountable BFT chain of $n$ validators tolerates up to a
$1/3$ adversarial fraction globally. Our adversary stays comfortably below that global bound, yet
buys enough faults to exceed the local threshold of a \emph{single} committee. The adversary is
thus globally sub-threshold but locally decisive. This is precisely the regime sharding creates
and monolithic analyses miss.

\paragraph{Why the cost of corruption collapses under sharding.}
Consider a network of $n$ validators partitioned into $m$ committees of size $N=n/m$. Attacking a
monolithic chain requires corrupting $\Theta(n)$ validators; attacking one shard requires
corrupting only $F+1 = \lfloor (N-1)/3\rfloor + 1$ validators. For the small committees used in
practice---RapidChain reports committee sizes on the order of a few tens of
nodes~\cite{zamani2018rapidchain}---the per-shard threshold is a handful of validators even when
$n$ is in the thousands. Random assignment and periodic reconfiguration raise the cost of
\emph{targeting} a specific committee, but they do not change the fact that, conditioned on
revealed membership, the local safety margin is small.

\paragraph{Why value is large.}
A shard's finalized state is acted upon as though it were permanent. Cross-shard transaction
protocols move assets between shards on the strength of finalized receipts, and bridges move
assets between chains the same way~\cite{zamyatin2021sok,lee2022sokbridge}. Extractable value $V$
is therefore bounded not by a single shard's local balances but by the liquidity reachable through
its cross-shard and bridge interfaces. The same plumbing that makes sharding useful makes a
single compromised shard valuable to attack.

\paragraph{Why timing is the crux.}
Accountable BFT does not prevent equivocation; it punishes it after the fact. The protective value
of slashing is therefore a function of \emph{latency}: how long after the conflicting certificates
appear can a transfer predicated on the bad branch still be reversed? Cross-shard pipelines are
frequently asynchronous: the source shard finalizes first, emits a receipt, and the destination
shard processes it later, paying gas on the destination side~\cite{harmony2019whitepaper}. If the
destination effect (and any subsequent withdrawal) clears before the source-side equivocation is
proven and the branch reverted, the adversary keeps the proceeds while the bribed validators bear
the slashing. This is the accountability race. Committee reconfiguration does not eliminate this
race; as we show in \S\ref{sec:reconf}, it \emph{relocates} it to the question of whether
old-epoch certified artifacts execute before old-epoch equivocation evidence is enforced.

\paragraph{Problem statement.}
Given a sharded BFT protocol with committee size $N$, a slashing/eviction mechanism with
accountability latency $\Tacc$, and a cross-shard settlement latency $\Tsettle$, we ask:
(i) under what conditions does there exist a window $\Tsettle < \Tacc$ in which fraudulently
finalized state can be converted to irreversible value; (ii) what is the minimum total bribe $B$
required to forge the necessary conflicting certificate; (iii) for which configurations is the
adversary's expected profit positive; and (iv) how does committee reconfiguration at the epoch
boundary change the answer? We seek conditions that are \emph{checkable} against a concrete
design, so that the same analysis serves both as an attack and as a design rule.
